\documentclass[11pt]{article}
\usepackage[margin=1in]{geometry}
\usepackage{amsmath,amssymb,amsthm}
\usepackage{array}
\usepackage{parskip}
\usepackage{booktabs}
\usepackage[fontset=fandol]{ctex}
\emergencystretch=2em

\newtheorem{theorem}{Theorem}
\newtheorem{lemma}[theorem]{Lemma}
\newtheorem{corollary}[theorem]{Corollary}
\newtheorem{proposition}[theorem]{Proposition}
\theoremstyle{definition}
\newtheorem{definition}[theorem]{Definition}
\theoremstyle{remark}
\newtheorem{remark}[theorem]{Remark}

\newcommand{\Z}{\mathbb{Z}}

\title{Refutation of conjecture \texttt{00000001109}}
\author{earthking11}
\date{2026-09-13}

\begin{document}
\maketitle

\section*{Verdict: FALSE}

We refute conjecture \texttt{00000001109} as stated, under the universal
reading that its own wording signals. For every \emph{irreducible} finite
non-crystallographic Coxeter group the degrees --- hence the exponents --- are
pairwise distinct, so every exponent has multiplicity exactly $1$ and the
conjectured conclusion (``some exponent of multiplicity $\ge 2$'') is false.

\medskip
\noindent\fbox{\parbox{0.94\textwidth}{\textbf{Scope (stated up front, and
honestly).} The refutation is \emph{universal}: it holds for all irreducible
finite non-crystallographic Coxeter groups, which is what the conjecture's
phrase ``universality of multiplicity structure'' (中文: 重数结构普遍性) says.
It is \emph{not} unconditional over arbitrary groups. Under an
\emph{existential} reading that also admits \emph{reducible}
non-crystallographic groups the conjecture is TRUE: the reducible group
$A_1 \times H_3$ is non-crystallographic and has exponent multiset
$\{1,1,5,9\}$, so the exponent $1$ has multiplicity $2$. See
Section~\ref{sec:scope}.}}

\section{The conjecture}

Quoted verbatim from \texttt{conjectures/00000001109.md}:

\begin{quote}
\textbf{English.} Definition: The exponents of a Coxeter system (the
characteristic polynomial). Conjecture: The exponent set of a
non-crystallographic Coxeter group contains, besides $(h, 1)$, some exponent of
multiplicity $\ge 2$ (universality of multiplicity structure).

\textbf{中文。} 定义：Coxeter 系统的指数(特征多项式)。猜想：非晶体 Coxeter
群的指数集包含 $(h, 1)$ 之外的重数 $\ge 2$ 的指数(重数结构普遍性)。
\end{quote}

The claim says: for a non-crystallographic Coxeter group, apart from the
special pair $(h,1)$, there is an exponent whose multiplicity is at least $2$.
The parenthetical ``universality of multiplicity structure'' indicates that this
is asserted as a general structural law across such groups, i.e.\ a
\emph{universal} statement over the class, not the mere existence of one
example. We take the conjecture in that universal sense.

\section{Conventions}

\begin{definition}[degrees, exponents, Coxeter number]
For a finite Coxeter group $W$ let $d_1 \le d_2 \le \dots \le d_n$ be its
\emph{degrees} (the degrees of the fundamental polynomial invariants) and let
$h$ be its \emph{Coxeter number}. The \emph{exponents} of $W$ are
\[
e_i \;=\; d_i - 1 \qquad (i = 1, \dots, n),
\]
and $h$ is the largest degree, $h = d_n$. The \emph{multiplicity} of an
exponent is its multiplicity as an element of the multiset
$\{e_1, \dots, e_n\}$.
\end{definition}

Since the exponents are the degrees minus one, distinctness of degrees is
equivalent to distinctness of exponents. For a finite Coxeter group the
exponents are exactly the roots of the characteristic polynomial; we use the
equivalent degree description throughout.

\begin{remark}[(h, 1) is vacuous under the standard convention]
\label{rem:hnotexp}
Because $h = d_n$ is the largest \emph{degree} and the exponents are the
degrees minus one, the largest exponent is $h - 1 < h$. Hence $h$ is
\emph{not} an exponent, and the phrase $(h,1)$ in the conjecture does not
correspond to an exponent at all: under the standard convention it is vacuous.
Even on the generous reading in which one adjoins $h$ to the exponent list, the
conclusion still fails (Proposition~\ref{prop:adjoinh} below).
\end{remark}

\section{The irreducible finite non-crystallographic types}

The complete list of irreducible finite Coxeter types that are
non-crystallographic is
\[
H_3, \qquad H_4, \qquad I_2(m)\ \ (m \ge 5).
\]
Here $I_2(3) = A_2$, $I_2(4) = B_2$, $I_2(6) = G_2$ are crystallographic, and
$I_2(2) = A_1 \times A_1$ is reducible; hence the irreducible
\emph{non-crystallographic} $I_2$ cases are exactly $m \ge 5$. (This
classification, and the exponent tables below, are classical; see
Humphreys~\cite{Humphreys} and Bourbaki~\cite{Bourbaki}.)

\begin{center}
\begin{tabular}{l l c l l}
\toprule
Type & Degrees & $h$ & Exponents & Multiplicities \\
\midrule
$H_3$ & $2,\,6,\,10$ & $10$ & $1,\,5,\,9$ & all $1$ \\
$H_4$ & $2,\,12,\,20,\,30$ & $30$ & $1,\,11,\,19,\,29$ & all $1$ \\
$I_2(m)$, $m \ge 5$ & $2,\,m$ & $m$ & $1,\,m-1$ & all $1$ \\
\bottomrule
\end{tabular}
\end{center}

In the last row the two exponents $1$ and $m-1$ coincide exactly when
$m = 2$; for $m \ge 5$ we have $m - 1 \ge 4$, so $1 \ne m - 1$. In every row
the exponents are pairwise distinct, so each occurs with multiplicity $1$.

\section{The argument}

\begin{theorem}\label{thm:main}
Let $W$ be an irreducible finite non-crystallographic Coxeter group. Then no
exponent of $W$ has multiplicity $\ge 2$.
\end{theorem}

\begin{proof}
By the classification of irreducible finite Coxeter types, $W$ is one of
$H_3$, $H_4$, or $I_2(m)$ with $m \ge 5$. From the exponent tables of
Section~3:
\begin{itemize}
  \item $H_3$: exponents $1, 5, 9$, which are pairwise distinct;
  \item $H_4$: exponents $1, 11, 19, 29$, pairwise distinct;
  \item $I_2(m)$, $m \ge 5$: exponents $1$ and $m-1$, distinct because
        $m - 1 \ge 4 > 1$.
\end{itemize}
In each case the exponents form a set with no repetitions, i.e.\ a multiset of
multiplicity $1$ everywhere. Therefore no exponent has multiplicity $\ge 2$.
\end{proof}

\begin{corollary}\label{cor:false}
Conjecture \texttt{00000001109}, read universally over irreducible finite
non-crystallographic Coxeter groups, is false.
\end{corollary}

\begin{proof}
Take $W = H_3$. Its exponent set is $\{1,5,9\}$; since $h = 10$, the pair
$(h,1)$ contributes nothing, and the remaining exponents $5, 9$ each have
multiplicity $1$. So the exponent set contains no exponent of multiplicity
$\ge 2$, contradicting the conjecture's conclusion. The same contradiction is
obtained from $H_4$ (exponents $1,11,19,29$) and from every $I_2(m)$ with
$m \ge 5$ (exponents $1, m-1$).
\end{proof}

\begin{proposition}\label{prop:adjoinh}
Adjoining $h$ to the exponent list does not create a repeated exponent:
\[
H_3:\ \{1,5,9,10\}, \qquad H_4:\ \{1,11,19,29,30\}
\]
are both all-distinct, and for $I_2(m)$ the list $\{1, m-1, m\}$ is
all-distinct for $m \ge 5$.
\end{proposition}

\begin{proof}
Immediate from the tables: $10$ differs from $1,5,9$; $30$ differs from
$1,11,19,29$; and for $m \ge 5$ the three numbers $1$, $m-1 \ge 4$, and $m$
are pairwise distinct (indeed $m-1 \ne m$).
\end{proof}

\section{Honest scope: universal versus existential reading}
\label{sec:scope}

The refutation above is a \emph{universal} statement over the irreducible
finite non-crystallographic types. We state clearly where the boundary of the
claim lies.

\begin{theorem}[Reducibility caveat]\label{thm:reducible}
Under an \emph{existential} reading that admits reducible non-crystallographic
Coxeter groups, conjecture \texttt{00000001109} is TRUE. Explicitly, the
reducible group $A_1 \times H_3$ is non-crystallographic and has exponent
multiset
\[
\{1\} \sqcup \{1, 5, 9\} \;=\; \{1, 1, 5, 9\},
\]
so the exponent $1$ has multiplicity $2 \ge 2$.
\end{theorem}

\begin{proof}
Exponents add over direct products of Coxeter groups: the exponents of
$W_1 \times W_2$ are the exponents of $W_1$ together with those of $W_2$.
The group $A_1$ has a single exponent, $1$ (degree $2$), and $H_3$ has
exponents $1, 5, 9$. Hence $A_1 \times H_3$ has exponent multiset
$\{1,1,5,9\}$, in which $1$ occurs twice. The direct product is
non-crystallographic because $H_3$ is, and it is not irreducible. Thus an
existential reading of the conjecture --- ``there exists a non-crystallographic
Coxeter group with a repeated exponent'' --- is satisfied.
\end{proof}

\begin{remark}[Why the universal reading is the intended one]
The conjecture does not say ``there exists a non-crystallographic Coxeter
group with \dots''; it says the exponent set \emph{contains} such an exponent
and labels this ``universality of multiplicity structure'' (重数结构普遍性).
``Universality'' is a statement about the class of objects, so the natural
reading is universal over non-crystallographic Coxeter groups. Under the
universal reading the conjecture is false, by
Corollary~\ref{cor:false}; the reducible counterexample
(Theorem~\ref{thm:reducible}) shows that the refutation depends on this
reading and would not survive an existential reformulation that admits
reducible groups. We prefer to state this openly rather than leave it implicit.
\end{remark}

\begin{remark}[On $(h,1)$]
As noted in Remark~\ref{rem:hnotexp}, under the standard convention $h$ is a
degree and not an exponent, so the phrase $(h,1)$ is vacuous. If one
nevertheless adjoins $h$ to the exponent list, the conclusion still fails
(Proposition~\ref{prop:adjoinh}). So the refutation does not hinge on the
interpretation of $(h,1)$: it fails under either reading, while the
reducibility issue above is the genuine scope boundary.
\end{remark}

\section{Reproducibility}

\paragraph{Python.} The script \texttt{reproduce.py} (Python~3 standard library
only) encodes the exponent lists of $H_3$, $H_4$, and $I_2(m)$ for
$m = 5, \dots, 12$, computes the multiplicities, asserts that every exponent has
multiplicity $1$, and asserts that the reducible case $A_1 \times H_3$ has a
repeated exponent (so the scope boundary is visible in code). It prints
\texttt{PASS} and exits $0$ exactly when every check holds:
\begin{quote}\ttfamily
\$ python3 reproduce.py\\
... (tables, multiplicity checks, scope check) ...\\
PASS: all checks verified; conjecture 00000001109 is FALSE ...
\end{quote}

\paragraph{Lean 4.} The accompanying project \texttt{lean4/} formalises the
finite multiplicity arithmetic in core Lean only (\texttt{import Std}, no
Mathlib, no \texttt{sorry}, no \texttt{native\_decide}). It proves, for the
concrete lists,
\[
\neg\,\exists x,\ x \in \{1,5,9\} \wedge 2 \le \#\{1,5,9\}(x)
\]
and the analogous statement for $\{1,11,19,29\}$ (axiom-free, by
\texttt{decide}), the general lemma that a duplicate-free list has every member
with count $1$, the parametric statement for $I_2(m)$, and --- openly --- the
reducibility boundary $\mathrm{count}_{1}(\{1,1,5,9\}) = 2$. It builds with
\texttt{lake build} and is audited by \texttt{lake env lean Check.lean}, which
reports no \texttt{sorryAx}. See \texttt{lean4/README.md} for the statement
table and the scope section.

\section{Cited inputs, not reproved}

The classification of irreducible finite Coxeter groups, the fact that the
exponents are the degrees minus one, and the exponent tables for
$H_3$, $H_4$, and $I_2(m)$ are classical. They are cited here, not reproved:
\begin{itemize}
  \item J.~E.~Humphreys, \emph{Reflection Groups and Coxeter Groups},
        Cambridge Studies in Advanced Mathematics 29, Cambridge University
        Press, 1990 (degrees, exponents, the classification).
  \item N.~Bourbaki, \emph{Groupes et algèbres de Lie}, Chapitres IV--VI,
        Hermann, Paris, 1968 (Coxeter groups and their invariant theory).
\end{itemize}
The present submission contributes the observation that the multiplicity
structure is trivial (all multiplicities $1$) for every irreducible finite
non-crystallographic type, hence the conjectured ``universality of multiplicity
structure'' fails; and it records explicitly the reducible exception that
bounds the claim.

\begin{thebibliography}{9}
\bibitem{Humphreys} J.~E.~Humphreys, \emph{Reflection Groups and Coxeter
Groups}, Cambridge Studies in Advanced Mathematics 29, Cambridge University
Press, 1990.
\bibitem{Bourbaki} N.~Bourbaki, \emph{Groupes et algèbres de Lie}, Chapitres
IV--VI, Hermann, Paris, 1968.
\end{thebibliography}

\end{document}
